\section[Text]{Text that behaves}

\blockframe{2}{Text that behaves}{Paragraphs, line breaks and lists}

\begin{frame}
    \frametitle{Where we are going}
    \vfill
    \begin{roadmap}{2}
        \roadmapresults
    \end{roadmap}
    \vfill
    \note{%
        The first block is done.
        It took 26 minutes.
        Say that out loud: it makes the remaining three blocks feel cheap.
    }
\end{frame}

\begin{frame}[fragile]
    \frametitle{Warm-up: three things are wrong here}
    \begin{columns}[T, onlytextwidth]
        \begin{column}{.55\textwidth}
            \tcbinputlst[fontsize=\footnotesize]{listings/02-warmup.tex}
        \end{column}
        \begin{column}{.4\textwidth}
            \begin{alertblock}<2->{Line 6}
                \code|\emph{| never closes
            \end{alertblock}
            \begin{alertblock}<2->{Line 9}
                A bare \code|%| eats the rest of the line
            \end{alertblock}
            \begin{alertblock}<2->{Line 10}
                \texttt{itemized} $\neq$ \texttt{itemize}
            \end{alertblock}
        \end{column}
    \end{columns}
    \note{%
        \textbf{Two minutes, laptops shut.}
        In pairs, out loud.
        \begin{itemize}
            \item Retrieval practice: pulling block 1 back out after a break
                is worth far more than you re-explaining it.
                It also wakes the room up.
            \item Do not reveal until at least two pairs have found all three.
                If the room is quiet, give them the first one.
        \end{itemize}
    }
\end{frame}

\begin{frame}
    \frametitle{Where do line breaks come from?}
    You type this:
    \tcbinputlst[][width=.6\textwidth]{listings/02-linebreaks.tex}

    \uncover<2->{%
        You get this:
        \begin{outputbox}[width=.9\textwidth]
            A freshly poured coffee is
too hot to drink, and a
forgotten one is too cold.

        \end{outputbox}
    }
    \vfill
    \begin{alertblock}<3->{}
        Three lines in, one line out
    \end{alertblock}
    \note{%
        \textbf{Predict first.}
        Ask the room what will come out before you click.
        Show of hands: three lines or one line?
        \begin{itemize}
            \item The wrong prediction is the point.
                Being surprised is what makes the rule stick.
                If you just state it, they nod and then write \texttt{\textbackslash\textbackslash} everywhere anyway.
        \end{itemize}
    }
\end{frame}

\begin{frame}[fragile]
    \frametitle{The rules, all of them}
    \begin{itemize}
        \item<1-> \textbf{Single newline}: a space
        \item<2-> \textbf{Empty line}: a new paragraph
        \item<3-> \code|%| at the end of a line: no space at all
        \item<4-> \code|~|: a space that never breaks, \code|Fig.~\ref{...}|
    \end{itemize}
    \vfill
    \begin{columns}[T, onlytextwidth]
        \begin{column}{.47\textwidth}
            \begin{exampleblock}<5->{One sentence per line}
                Same PDF, cleaner \texttt{git diff}
            \end{exampleblock}
        \end{column}
        \begin{column}{.47\textwidth}
            \begin{alertblock}<6->{\texttt{\textbackslash\textbackslash} is \textbf{not} a paragraph break}
                Only in tables and maths
            \end{alertblock}
        \end{column}
    \end{columns}
    \note{%
        The \texttt{\textbackslash\textbackslash} warning is the most valuable thing in block 2.
        \begin{itemize}
            \item \texttt{\%} at the end of a line is useful in tables and macros.
            \item One sentence per line: \texttt{git diff} then shows the sentence that changed, not the paragraph.
                The template is written this way.
            \item \texttt{\textbackslash\textbackslash} ends a line; anywhere but tables and maths it gives underfull boxes.
            \item Every beginner reaches for it because it feels like Word's Enter key.
                It is wrong in a way that only shows once the document is long enough to reflow.
            \item Name the symptom, so they recognise it later:
                \texttt{Underfull \textbackslash hbox (badness 10000)} in the log,
                and a paragraph whose last line is stretched across the full width.
        \end{itemize}
    }
\end{frame}

\begin{frame}
    \frametitle{Paragraph spacing: do it once}
    \begin{columns}[T, onlytextwidth]
        \begin{column}{.47\textwidth}
            {\large\bfseries What beginners do\par}
            \medskip
            \tcbinputlst{listings/02-noindent.tex}
            \begin{alertblock}{}
                Every paragraph, every time
            \end{alertblock}
        \end{column}
        \begin{column}{.47\textwidth}
            {\large\bfseries What you do\par}
            \medskip
            \tcbinputlst{listings/02-parskip.tex}
            \begin{exampleblock}{}
                Once, in the preamble
            \end{exampleblock}
        \end{column}
    \end{columns}
    \note{%
        \begin{itemize}
            \item Miss one \texttt{\textbackslash noindent} and it looks broken;
                undo \texttt{parskip} by deleting one line.
            \item If you find yourself repeating a command, look for a package.
            \item Structure before appearance, for the third time.
                By now they should feel the pattern rather than need it explained.
            \item A blank line between paragraphs is the normal convention in German-language typesetting anyway,
                so for a lot of their coursework \texttt{parskip} is simply correct.
        \end{itemize}
    }
\end{frame}

% Each list is typeset twice from the same file: once as code, once for real.
\begin{frame}
    \frametitle{Three kinds of list}
    \begin{columns}[T, onlytextwidth]
        \begin{column}{.31\textwidth}
            \tcbinputlst[fontsize=\footnotesize]{listings/02-itemize.tex}
            \begin{outputbox}[height=1.6cm]
                \begin{itemize}
    \item bread
    \item milk
\end{itemize}

            \end{outputbox}
            \textbf{itemize}: order does not matter
        \end{column}
        \begin{column}{.31\textwidth}
            \tcbinputlst[fontsize=\footnotesize]{listings/02-enumerate.tex}
            \begin{outputbox}[height=1.6cm]
                \begin{enumerate}
    \item heat water
    \item pour
\end{enumerate}

            \end{outputbox}
            \textbf{enumerate}: order \emph{does} matter
        \end{column}
        \begin{column}{.31\textwidth}
            \tcbinputlst[fontsize=\footnotesize]{listings/02-description.tex}
            \begin{outputbox}[height=1.6cm]
                \begin{description}
    \item[mug] ceramic
    \item[flask] steel
\end{description}

            \end{outputbox}
            \textbf{description}: naming things
        \end{column}
    \end{columns}
    \vfill
    \begin{exampleblock}<2->{}
        Same shape: \texttt{\textbackslash begin}, \texttt{\textbackslash item}s, \texttt{\textbackslash end}
    \end{exampleblock}
    \note{%
        Fast, three minutes.
        They already know what a list is; they need the three names and the shared shape.
        \begin{itemize}
            \item Lists nest, and the labels sort themselves out.
                Worth demoing live.
                The automatic change of label (bullet to dash, 1 to a) is another small
                ``it did that for me'' moment.
            \item \texttt{description} is the one they have never seen.
                The bracket argument is the term being defined.
                TODO 2.6 uses it for the symbol list, which becomes maths in block 3.
        \end{itemize}
    }
\end{frame}

\begin{frame}
    \frametitle{The preamble: where the settings live}
    \only<1>{\tcbinputlst[highlightlines={1}]{listings/02-preamble.tex}}%
    \only<2>{\tcbinputlst[highlightlines={3-6}]{listings/02-preamble.tex}}%
    \only<3->{\tcbinputlst[highlightlines={8}]{listings/02-preamble.tex}}%
    \begin{itemize}
        \item<4-> Between \texttt{\textbackslash documentclass} and \texttt{\textbackslash begin\{document\}}
        \item<5-> \texttt{\textbackslash usepackage}: load an extension (CTAN: about 6\,000)
        \item<6-> \texttt{hyperref} goes last
    \end{itemize}
    \note{%
        \begin{itemize}
            \item ``Can \LaTeX{} do\ldots?'' Almost always, with one line in the preamble.
            \item Package order rarely matters; \texttt{hyperref} is the one to remember.
            \item Held back until now on purpose.
                In block 1 the preamble was ``the settings area, we come back to it'',
                which was enough because they had nothing to put in it yet.
                Now \texttt{amsmath} is waiting for the next block.
            \item 6\,000 always gets a reaction.
                The follow-up: you will personally use about eight of them.
        \end{itemize}
    }
\end{frame}

\begin{frame}
    \frametitle{Your turn}
    \handson{2}{2.1--2.6}{5--10}{checkpoint-02.tex}{%
        Work down the file.\par
        \medskip
        {\normalsize\color{black!60}
        add \texttt{parskip} $\rightarrow$ write a second paragraph
        $\rightarrow$ fix the three line-break crimes in Setup
        $\rightarrow$ an \texttt{itemize}, an \texttt{enumerate} and a \texttt{description}}
    }
    \note{%
        \textbf{5 to 10 minutes.}
        \begin{itemize}
            \item TODO 2.3 is the interesting one: the paragraph has the \texttt{\textbackslash\textbackslash} habits baked in,
                and it compiles perfectly happily.
                \LaTeX{} will not tell you; you have to look at the page.
            \item Watch for a blank line \emph{inside} a list and strange spacing.
                Common, harmless, worth a word from the front if you see it twice.
            \item \textbf{Break after this}, 5 minutes.
                Again: name the restart time.
        \end{itemize}
    }
\end{frame}
